% 英文摘要
\chapter*{\hfill ABSTRACT \hfill}
\addcontentsline{toc}{chapter}{ABSTRACT}

\par SCADA (Supervisory Control and Data Acquisition) systems serve as the central
nervous system of critical national infrastructure, widely deployed in gas, power,
and water industrial control scenarios. Security breaches directly threaten public
safety and economic stability. With the deep integration of Industrial Internet of
Things (IIoT) and artificial intelligence, traditional signature-based intrusion
detection methods can no longer cope with increasingly sophisticated cyber attacks.
There is an urgent need to develop intelligent, lightweight intrusion detection
methods tailored for SCADA systems.

\par However, SCADA scenarios impose unique and stringent requirements on detection
models. On one hand, the temporal structure of industrial control protocols (e.g.,
Modbus) contains rich domain semantics that general time-series models fail to
exploit. On the other hand, constrained by the limited computing power and storage
of edge devices, detection models must balance accuracy with deployability. Existing
work often re-uses generic network intrusion detection features without protocol-level
semantic encoding, and few studies systematically explore the complete optimization
chain from algorithm design to on-chip deployment.

\par To address these issues, this thesis takes the IanArffDataset (274{,}000
SCADA gas-pipeline records, 20 raw features, three-layer labels) as the research
object, and conducts a full-stack study covering data features, model architecture,
and edge deployment, all centered on the core goal of ``lightweight intrusion
detection''. The main contributions are as follows:

\par (1) A domain feature engineering method based on SCADA protocol semantics is
proposed. By analyzing the 4-row master-slave interaction structure of Modbus
and treating 4 rows as one Modbus cycle, 27-dimensional protocol semantic features
(3 row-level + 8 window-level) are designed on top of the 17-dimensional minimal
feature set. Comparison experiments on multiple models (LightGBM, TCN, etc.)
demonstrate that the proposed feature set consistently outperforms the original
44-dimensional features, validating the conclusion that ``the marginal benefit of
domain feature engineering exceeds that of model upgrading''.

\par (2) A lightweight intrusion detection model (TCN-SE) that integrates SE
attention into Temporal Convolutional Network is proposed. Through a systematic
comparison of 27 typical time-series models, the inductive bias of TCN's causal
convolution is found to naturally fit Modbus short-term temporal structure.
Further introducing the SE channel attention, with only 3{,}300 additional parameters
(+4.3\%), the model achieves Macro-F1 = 0.8340 and PR-AUC = 0.9025 on the test
set, marking the first time a deep learning model surpasses gradient boosting
trees (LightGBM, Macro-F1 = 0.8337) on this task. Ablation studies on window
length ($w=4, 8, 16$) and channel width (ch=8 $\sim$ 64) reveal that ``8-step
window (2 Modbus cycles) is the Pareto-optimal structure for this dataset'' and
verify the role of SE in improving minority-class recall.

\par (3) An extreme compression and mixed-precision quantization scheme for MCU
deployment is proposed. Through systematic channel width ablation, a TCN V19
variant with only 4{,}237 parameters is obtained, and the PR-AUC counterintuitively
rises to 0.9133 (+0.67\%), challenging the conventional belief that ``larger models
yield better performance''. A ``weight-INT8 + activation-FP32'' hybrid quantization
scheme is further proposed, achieving a 5.5\,KB model with 0.47\,ms inference latency
on the Cortex-M4 168\,MHz platform. The prediction consistency before and after
quantization reaches 99.80\%, with SNR $\geq$ 38\,dB, indicating essentially lossless
quantization. A complete 380-line C99 inference implementation is open-sourced and
can be directly deployed on STM32, ESP32, Arduino and other embedded platforms.

\par Following a progressive ``data-algorithm-chip'' research route, this thesis
forms a complete innovation chain from domain features to MCU deployment. The
three innovations achieve the best Macro-F1 = 0.8340 (single model) and
PR-AUC = 0.9133 (compressed single model) on the public dataset, and verify the
feasibility of maintaining industrial-grade detection performance on a 5.5\,KB
embedded model, providing a reproducible reference for the engineering deployment
of SCADA intrusion detection.

\vspace{1em}

\noindent\textbf{Keywords:} SCADA; Intrusion Detection; Temporal Convolutional
Network; Model Compression; Mixed-Precision Quantization; Embedded Deployment

\clearpage
